\chapter{Data Construction Details}
\label{app:data}

\section{Data Audit}

Table~\ref{tab:audit_full} reports the raw and cleaned row counts and the main
data defects. The trailing-period duplicates are normalized before deduplication;
the group-like targets are removed from the candidate table and kept in the group
corpus. The cleaned table retains the original row index so that every row can be
traced back to the raw ANES table.

\begin{table}[h]
\centering
\small
\begin{tabular}{lr}
\toprule
Item & Count \\
\midrule
Raw rows & 132,948 \\
Trailing-period duplicate rows & 8,165 \\
Group targets mixed into candidates & 23,060 \\
Missing or out-of-range scores & 3,588 \\
Duplicate rows removed after normalization & 8,141 \\
Clean rows & 98,159 \\
\bottomrule
\end{tabular}
\caption{Raw and cleaned row counts.}
\label{tab:audit_full}
\end{table}

\section{Splits}

\begin{table}[h]
\centering
\small
\begin{tabular}{llrr}
\toprule
Split & Years & Rows & Respondents \\
\midrule
\texttt{train\_t} & 2016--2020 & 6,516 & all \\
\texttt{val\_t} & 2022 & 2,501 & sampled \\
\texttt{test\_t} & 2024 & 3,421 & 1,712 \\
\texttt{train\_r} & 2016--2020 & 5,094 & 80\% \\
\texttt{test\_r} & 2016--2022 & 6,014 & 192 \\
\texttt{test\_tr} & 2024 & 226 & 113 \\
\bottomrule
\end{tabular}
\caption{Candidate splits. The two protocols are alternative designs and overlap
across each other; within each protocol the splits are disjoint.}
\label{tab:splits}
\end{table}

\section{Sampling Strategies}

The training pool is drawn with nine complementary strategies: balanced coverage
of the score range, hard cases near the middle of the scale, party-by-ideology
cells, similarity-stratified bins, cross-pressure cases, extreme scores,
incumbent effects, temporal coverage, and candidate-specific balance. The pool
contains 6,516 rows and is intentionally not a uniform sample; the test sets are
built from all cleaned rows in the target years so that evaluation is not shaped
by the same sampling design.

\section{Feature Definitions}

The core features are party match (+1 same party, $-1$ independent, 0 otherwise),
seven-point ideological distance, a composite demographic similarity index, and
three cross-pressure indicators: CP1 (same party, ideological distance $\ge3$),
CP2 (different party, distance $\le1$), and CP3 (similarity $\ge0.7$, different
party).
